% Ad Atticum 2.16 (ad-atticum-02-16) --- English translation
% Translated by Claude (Anthropic) from the Latin (Perseus).
% Style guide: STYLE.md.

\ciceroLetterOpener{Cicero to Atticus, greetings}

\ciceroSection{1} On the day before the Kalends of May, when I had dined and was already going to sleep, that letter was delivered to me in which you write about the Campanian land. What more? At first it so stung me as to take away my sleep, but more by reflection than by trouble. As I thought it over, the following considerations occurred to me. First: from what you had written in earlier letters --- that you had heard from one of his intimates that something would be brought forward which no one would disapprove --- I had feared something larger than this. This did not appear to me of that kind. Next, by way of consoling myself: every expectation of agrarian generosity seems to have been turned aside on to the Campanian land --- which, even at ten iugera each, cannot support more than five thousand men; the rest of the multitude must be alienated from those leaders. Further, if there is any one thing that can more vehemently inflame the spirits of the good, whom I see already stirred up, it is this --- and the more so, since with the tolls of Italy abolished and the Campanian land divided up, what domestic revenue is left save the twentieth? which seems to me bound to perish in one little public meeting at the shouts of the slave-attendants of the leaders.

\ciceroSection{2} Our Gnaeus indeed I no longer plainly know what he is thinking. ``For he blows now no longer with little pipes / but with savage bellows without the muzzle.'' --- that he could even be brought to that point! For up to now he was sophisticating thus: he approved Caesar's laws; for the actions Caesar himself must answer; the agrarian law had pleased him --- whether it could have been vetoed or not did not concern him; concerning the king of Alexandria, he was pleased that the matter be settled at last; whether Bibulus was watching the sky at the time or not was not his to investigate; about the publicans, he wished to oblige that order; what would have happened if Bibulus had then come down to the Forum he could not divine. But now, Sampsiceramus --- what will you say? That you have set up a tribute for us on Mount Antilibanus, and have taken away the Campanian land? Well? How will you maintain this? ``I shall hold you down,'' he says, ``crushed, by Caesar's army.'' But, by Hercules, you hold me down not so much by that army as by the ungrateful spirits of the men who are called good, who have never returned to me any fruit or thanks, not only of rewards but not even of mere conversations.

\ciceroSection{3} If I were to rouse myself to that side, I should certainly find some way of resisting. Now I have plainly settled this: since there is such a great quarrel between your friend Dicaearchus and my friend Theophrastus, that yours puts the active life \textit{[Greek: ton praktikon bion]} far before everything else, and this one the contemplative \textit{[Greek: ton the\=or\=etikon]} --- I shall be seen to have served the wishes of both. For I think I have given Dicaearchus full satisfaction; I look back now to that other family, which not only allows me to rest but rebukes me because I have not always been quiet. Wherefore let us bend, our dear Titus, to those splendid pursuits, and at last go back to the place from which we ought never to have departed.

\ciceroSection{4} What you write about Quintus's letter --- to me also it has been ``a lion in front, a snake behind \ldots'' \textit{[Greek: prosthe le\=on, opithen de ---]} (Iliad 6.181). What to say I do not know. For he so laments his stay in the first verses that he could move anyone; then again he so loosens up that he asks me to correct and edit his annals. That which you write, however, I should like you to attend to: about the toll on the carriage. He says that, on the advice of the council, he has referred the matter to the senate. He had not yet, evidently, read my letter, in which I had written, after considering and exploring the matter, that nothing was due. If any Greeks have now come to Rome from Asia on this case, please see them; and if it seems good to you, show them what I think of the matter. If I can leave off, lest the best cause perish in the senate, I shall give the publicans satisfaction. If not (I shall speak with you frankly), in this case I am rather for all Asia and its businessmen; for it concerns them vehemently too. This much I feel: that the matter greatly concerns us. But you will see to it. As for the quaestors, please, are they doubting even about the cistophorus? For if there is nothing else, when I have tried all, I shall not despise even that last resort: we shall see you at the Arpinum estate and shall receive you with country hospitality, since this seaside hospitality you have despised.
